\documentclass[10pt,a4paper]{amsart}
\usepackage[margin=19mm]{geometry}
\usepackage{amsmath,amssymb,mathtools}
\usepackage[scr=rsfs,cal=euler]{mathalfa}
\usepackage{xfrac}
\usepackage[unicode,hidelinks,bookmarks=false]{hyperref}
\newcommand{\ie}{i.e.\ }
\newcommand{\cf}{cf.\ }
\newcommand{\resp}{respectively\ }
\newcommand{\Subsection}{\subsection}
\providecommand{\nobreakdash}{\nobreak\hbox{-}}
\setlength{\emergencystretch}{2em}
\multlinegap0pt
\newcommand{\C}{\mathbb{C}}
\newcommand{\Q}{\mathbb{Q}}
\newcommand{\T}{\mathbb{T}}
\newcommand{\Var}{\operatorname{Var}}
\newcommand{\Z}{\mathbb{Z}}
\newcommand{\dd}{\,\mathrm{d}}
\newcommand{\wind}{\operatorname{wind}}
\theoremstyle{plain}
\newtheorem{theorem}{Theorem}[section]
\newtheorem{proposition}[theorem]{Proposition}
\newtheorem{lemma}[theorem]{Lemma}
\newtheorem{corollary}[theorem]{Corollary}
\theoremstyle{remark}
\newtheorem{remark}[theorem]{Remark}
\title[The critical mixed-arithmetic four-point HRT theorem]{The critical mixed-arithmetic four-point HRT theorem}
\author{Vignon Oussa}
\date{Source: arXiv:2609.27970v1 (CC BY 4.0)}
\begin{document}
\maketitle

\noindent\emph{Source and licence.} The delivered work is the article shown in the title above, version arXiv:2609.27970v1 (CC BY 4.0), distributed under the Creative Commons Attribution 4.0 International licence, as recorded in the arXiv licence metadata for this exact version and shown on the abstract page saved with this item. The full licence notice, which carries the licence URL and the address of that abstract page, is the bundled file \texttt{licenses/arxiv-2609.27970-CC-BY-4.0.txt}.
\par\smallskip

\noindent\textit{Editorial scope.} Counted unit: the article's Lemma (source label lem:winding) (source0.tex lines 1067--1076, source label \texttt{lem:winding}): ``Measurable winding obstruction Let ( ), let (0 h L 2( ) ), and let (b C 2( ; ) ). If h(x- ) b(x)h(x) almost everywhere, then ( (b) 0 ).''. It is delivered together with the article's complete original proof (source0.tex lines 1078--1177, one proof environment adjacent to the statement, 100 lines), copied line for line from the article's own text. Required context retained uncounted: eq:cf-return (lines 1034--1037). Editorial formatting only: an AMS \texttt{amsart} preamble that restates the article's own macro definitions and its own theorem declarations, and the theorem environments the retained lines use; the article's own class file is neither shipped nor input; the retained environments are renumbered sequentially in order of appearance, whereas the article prints them with its own numbering; the article's equation numbering is not reproduced and every cross-reference inside the retained text resolves to the wrapper's numbering. No citation occurs inside any retained line, so the wrapper carries no bibliography; All retained bodies are exact copies of the bundled \texttt{sources/211538/source0.tex}, so no omission receipt applies. No asset-inclusion command occurs inside any retained span and the package ships no image asset.


\section*{Retained context: eq:cf-return (source0.tex lines 1034--1037)}

\begin{equation}
  q_n\theta-p_n\longrightarrow0.
  \label{eq:cf-return}
\end{equation}


\section{The counted unit: Lemma (source label lem:winding) and its complete original proof}

\begin{lemma}[Measurable winding obstruction]
\label{lem:winding}
Let \(\theta\notin\Q\), let \(0\neq h\in L^2(\T)\), and let
\(b\in C^2(\T;\C^\times)\). If
\begin{equation}
  h(x-\theta)=b(x)h(x)
  \label{eq:abstract-cocycle}
\end{equation}
almost everywhere, then \(\wind(b)=0\).
\end{lemma}

\begin{proof}
\textbf{Step 1: Remove the zero set of the transfer function.}
Because \(b\) never vanishes, the zero set of \(h\) is invariant, modulo null
sets, under the irrational rotation. Ergodicity implies that this zero set has
measure zero or one. Since \(h\neq0\), it has measure zero. Therefore
\[
  u(x)=\frac{h(x)}{|h(x)|}
\]
is defined almost everywhere and has modulus one.

\smallskip
\noindent
\textbf{Step 2: Separate the winding from the periodic phase.}
Let \(\psi=b/|b|\). Then
\[
  u(x-\theta)=\psi(x)u(x)
\]
almost everywhere and \(\wind(\psi)=\wind(b)\). If this common winding is
\(d\in\Z\), there is a real-valued periodic \(C^2\) function \(g\) such that
\[
  \psi(x)=e^{2\pi i(dx+g(x))}.
\]

\smallskip
\noindent
\textbf{Step 3: Iterate through continued-fraction returns.}
Iteration through \(q_n\) steps gives
\[
  u(x-q_n\theta)=e^{2\pi i\Phi_n(x)}u(x),
\]
where
\begin{equation}
  \Phi_n(x)
  =dq_nx-d\theta\frac{q_n(q_n-1)}2
   +\sum_{j=0}^{q_n-1}g(x-j\theta).
  \label{eq:Phi-def}
\end{equation}
Since \(|u|=1\),
\[
  |e^{2\pi i\Phi_n(x)}-1|
  =|u(x-q_n\theta)-u(x)|.
\]
By \eqref{eq:cf-return} and translation continuity in \(L^2(\T)\),
\begin{equation}
  e^{2\pi i\Phi_n}\longrightarrow1
  \quad\text{in }L^2(\T).
  \label{eq:phase-to-one}
\end{equation}
In particular,
\begin{equation}
  \int_\T e^{2\pi i\Phi_n(x)}\dd x\longrightarrow1.
  \label{eq:phase-integral-one}
\end{equation}

\smallskip
\noindent
\textbf{Step 4: Exclude nonzero winding by oscillation.}
Assume \(d\neq0\). Differentiating \eqref{eq:Phi-def} gives
\[
  \Phi_n'(x)
  =dq_n+\sum_{j=0}^{q_n-1}g'(x-j\theta).
\]
The function \(g'\) has mean zero and bounded variation. Denjoy--Koksma
therefore yields
\[
  \sup_x\left|
    \sum_{j=0}^{q_n-1}g'(x-j\theta)
  \right|
  \leq\Var(g').
\]
Hence \(\inf_x|\Phi_n'(x)|\geq c q_n\) for all large \(n\), with \(c>0\).
Moreover,
\[
  \|\Phi_n''\|_{L^1}
  \leq q_n\|g''\|_{L^1}.
\]

Although \(\Phi_n\) need not be periodic, it satisfies
\[
  \Phi_n(x+1)=\Phi_n(x)+dq_n,
\]
and \(\Phi_n'\) is periodic. Consequently, the boundary term in integration
by parts vanishes:
\[
  \left[
    \frac{e^{2\pi i\Phi_n(x)}}{2\pi i\Phi_n'(x)}
  \right]_{0}^{1}=0.
\]
It follows that
\[
  \left|
    \int_\T e^{2\pi i\Phi_n(x)}\dd x
  \right|
  \leq
  \frac{\|\Phi_n''\|_{L^1}}
       {2\pi(\inf|\Phi_n'|)^2}
  =O(q_n^{-1}).
\]
This contradicts \eqref{eq:phase-integral-one}. Therefore \(d=0\).
\end{proof}

\end{document}
